\documentclass[12pt, twoside]{article}
%\documentclass[12pt, twoside]{article}
\usepackage[letterpaper, margin=1in, headsep=0.2in]{geometry}
\setlength{\headheight}{0.6in}
%\usepackage[english]{babel}
\usepackage[utf8]{inputenc}
\usepackage{microtype}
\usepackage{amsmath}
\usepackage{amssymb}
%\usepackage{amsfonts}
\usepackage[nomessages]{fp} %\FPeval{\var-name}{2*sin(pi/6)}
\usepackage{siunitx} %units in math. eg 20\milli\meter
\usepackage{yhmath} % for arcs, overparenth command
\usepackage{tikz} %graphics
\usetikzlibrary{quotes, angles, arrows, arrows.meta}
\usepackage{pgfplots}
\usepgfplotslibrary{fillbetween}
\pgfplotsset{compat=1.16}
\usepackage{graphicx} %consider setting \graphicspath{{images/}}
\usepackage{parskip} %no paragraph indent
\usepackage{enumitem}
\usepackage{multicol}
\usepackage{venndiagram}

\usepackage{fancyhdr}
\pagestyle{fancy}
\fancyhf{}
\renewcommand{\headrulewidth}{0pt} % disable the underline of the header
\raggedbottom
\hfuzz=2mm %suppresses overfull box warnings

\usepackage{hyperref}

\fancyhead[LO]{La Scuola d'Italia / Dr.\ Huson / IB Math 11 \\* 21 October 2026}
\fancyhead[RO]{\fbox{\begin{minipage}{6cm}\footnotesize Nickname:\\[0.5cm]\end{minipage}}}
\fancyfoot[C]{\thepage}

\begin{document}
\subsubsection*{2.3 Pre-Quiz: Sequences, Lines, Radicals and Exponents}

\begin{enumerate}[itemsep=0.15cm]

%--- Problem 1: Geometric sequence 24, 12, 6: r, u_6, why S_inf exists; find r from u_1 and S_inf [6 marks] ---
%--- Source: Original (freshly authored, AA SL 1.2, 1.3); review of 1-9PQ, 1-11CW ---
\item[1.]
Consider the geometric sequence
$$24,\; 12,\; 6,\; \ldots$$

\begin{enumerate}[itemsep=0.1cm]
  \item Write down the common ratio $r$.
  \item Find $u_{6}$.
  \item Explain why the series $24 + 12 + 6 + \cdots$ has a sum to infinity.
\end{enumerate}

A different geometric series has first term $10$ and sum to infinity $15$.

\begin{enumerate}[itemsep=0.1cm]
  \item[(d)] Find the common ratio of this series.
\end{enumerate}

\begin{tikzpicture}
  \draw (-0.6,0) rectangle (14.6,15.4);
  \foreach \y in {1,2,...,15} \draw[dotted] (0,\y)--(14.1,\y);
\end{tikzpicture}

\newpage

%--- Problem 2: Points P, Q; midpoint, gradient of PQ, perpendicular bisector L, point (k, 9) on L [9 marks] ---
%--- Source: Original, modelled on AA SL 2022 May P1 TZ1 Q1 and AA SL 2023 May P1 TZ1 Q1 (AA SL 2.1) ---
\item[2.]
Point $P$ has coordinates $(-4,\ 1)$, and point $Q$ has coordinates
$(6,\ -3)$. Point $M$ is the midpoint of $[PQ]$.

\begin{enumerate}[itemsep=0.1cm]
  \item Find the coordinates of $M$.
  \item Find the gradient of $[PQ]$.
\end{enumerate}

The line $L$ is perpendicular to $[PQ]$ and passes through $M$.

\begin{enumerate}[itemsep=0.1cm]
  \item[(c)] Write down the gradient of $L$.
  \item[(d)] Find the equation of $L$.
  \item[(e)] The line $L$ passes through the point $(k,\ 9)$. Find the value
  of $k$.
\end{enumerate}

\begin{tikzpicture}
  \draw (-0.6,0) rectangle (14.6,15.4);
  \foreach \y in {1,2,...,15} \draw[dotted] (0,\y)--(14.1,\y);
\end{tikzpicture}

\newpage

%--- Problem 3: f(x) = 2x - 3; intercepts; parallel g through origin; perpendicular h through (2, 6); intersection of f and h; x-intercept of h [12 marks] ---
%--- Source: Original, modelled on AA SL 2022 Nov P1 Q1 (AA SL 2.1) ---
\item[3.]
Let $f(x) = 2x - 3$, for $x \in \mathbb{R}$.

\begin{enumerate}[itemsep=0.1cm]
  \item Write down the $y$-intercept of the graph of $f$, and find the
  $x$-intercept.
  \item The graph of a linear function $g$ is parallel to the graph of $f$
  and passes through the origin. Find an expression for $g(x)$.
  \item The graph of a linear function $h$ is perpendicular to the graph of
  $f$ and passes through the point $(2,\ 6)$. Find an expression for $h(x)$.
  \item Find the coordinates of the point where the graphs of $f$ and $h$
  intersect.
  \item The graph of $h$ meets the $x$-axis at the point $R$. Find the
  coordinates of $R$.
\end{enumerate}

\begin{tikzpicture}
  \draw (-0.6,0) rectangle (14.6,14.4);
  \foreach \y in {1,2,...,14} \draw[dotted] (0,\y)--(14.1,\y);
\end{tikzpicture}

\newpage

%--- Problem 4: 4^x * 2^(x+1) = 2^(ax+b); solve = 16; solve 9^x = 27; show (sqrt12 + sqrt3)^2 = 27; 6/(3 - sqrt3) in the form a + b sqrt3 [13 marks] ---
%--- Source: Original, in the style of AA SL Paper 1 prior-learning questions (AA SL 1.5, 1.7) ---
\item[4.]

\begin{enumerate}[itemsep=0.1cm]
  \item The expression $4^{x} \times 2^{x + 1}$ can be written in the form
  $2^{ax + b}$. Find the value of $a$ and the value of $b$.
  \item Hence, solve the equation $4^{x} \times 2^{x + 1} = 16$.
  \item Solve the equation $9^{x} = 27$.
  \item Show that $\left(\sqrt{12} + \sqrt{3}\right)^{2} = 27$.
  \item Write $\dfrac{6}{3 - \sqrt{3}}$ in the form $a + b\sqrt{3}$,
  where $a, b \in \mathbb{Z}$.
\end{enumerate}

\begin{tikzpicture}
  \draw (-0.6,0) rectangle (14.6,15.4);
  \foreach \y in {1,2,...,15} \draw[dotted] (0,\y)--(14.1,\y);
\end{tikzpicture}

\end{enumerate}
\end{document}
